\documentclass[11pt,a4paper]{article}
\usepackage[margin=2.5cm]{geometry}
\usepackage{amsmath,amssymb,amsthm}
\usepackage[hidelinks]{hyperref}
\usepackage{xurl}
\usepackage{booktabs}
\usepackage{microtype}
\usepackage{tikz-cd}

\newtheorem{theorem}{Theorem}
\newtheorem{proposition}[theorem]{Proposition}
\newtheorem{corollary}[theorem]{Corollary}
\theoremstyle{definition}
\newtheorem{definition}[theorem]{Definition}
\theoremstyle{remark}
\newtheorem{remark}[theorem]{Remark}

\newcommand{\Z}{\mathbb{Z}}
\newcommand{\Q}{\mathbb{Q}}
\newcommand{\C}{\mathbb{C}}
\newcommand{\mfld}[1]{\texttt{#1}}
\newcommand{\Gal}{\mathrm{Gal}}
\newcommand{\SU}{\mathrm{SU}}
\newcommand{\Weyl}{\mathrm{Weyl}}
\newcommand{\disc}{\mathrm{disc}}
\newcommand{\vol}{\mathrm{vol}}

\begin{document}

\title{The Meyerhoff Manifold as a Dual Dehn Filling:\\
Arithmetic Independence of $\SU(2)$ and $\SU(4)$ Cusped Parents}

\author{Marvin L.~Gentry%
\thanks{Independent Researcher, Seattle, WA, USA;
\texttt{drmlgentry@protonmail.com};
ORCID: 0009-0006-4550-2663}}
\date{\today}
\maketitle

\begin{abstract}
The Meyerhoff manifold $M_{\mathrm{PMNS}}$ -- the minimum-volume
closed hyperbolic 3-manifold with first homology $\Z/5$ -- admits
two distinct Dehn filling descriptions from arithmetically
independent cusped parents: $M_{\mathrm{PMNS}} = \mfld{m003}(-2,3)
= \mfld{m019}(2,1)$. The cusp field of \mfld{m003} has Galois group
$\Z/2$; the cusp field of \mfld{m019} has Galois group $S_4$. We
prove these two cusp fields are linearly disjoint over $\Q$, and
that the Galois group of the Galois closure of their compositum is
$S_4\times\Z/2$, order 48. We further prove that the cusp field
of \mfld{m003} is linearly disjoint from the cusp field of
\mfld{m006} (the parent of the CKM manifold), yielding a Galois
closure with group $D_6$, order 12 -- not a direct product of
Weyl groups. We report a negative result: a systematic search over
9 candidate manifolds with $S_4$-type cusp fields, at Dehn surgery
slopes $|p|,|q|\le 15$, finds no dual filling description of the
companion manifold $M_{\mathrm{CKM}}$. We note, without asserting
a physical derivation, that the group $S_4\times\Z/2$ is
abstractly isomorphic to $\Weyl(\SU(4))\times\Weyl(\SU(2))$.
All computations are verified by SnapPy and Sage at 300-bit
precision.
\end{abstract}

\section{Introduction}

The Hyperbolic Flavor Geometry (HFG) program associates compact
hyperbolic 3-manifolds to Standard Model flavor parameters. The
lepton mixing matrix (PMNS) is encoded by the closed manifold
$M_{\mathrm{PMNS}}$, obtained by Dehn filling a cusped parent
manifold.

By Mostow rigidity, every complete finite-volume hyperbolic
structure on a 3-manifold of dimension $\ge 3$ is unique up to
isometry; geometric invariants such as volume and cusp shape are
therefore topological invariants of the manifold, not artifacts of
a choice of metric.

This paper proves the following results.

\begin{theorem}[Dual surgery]\label{thm:dual}
$M_{\mathrm{PMNS}} = \mfld{m003}(-2,3) = \mfld{m019}(2,1)$: the
Meyerhoff manifold arises as a Dehn filling of two distinct cusped
parents, \mfld{m003} and \mfld{m019}.
\end{theorem}

\begin{theorem}[Linear disjointness of $K_{m003}$ and $K_{m019}$]\label{thm:disjoint}
The cusp fields $K_{\mfld{m003}}$ and $K_{\mfld{m019}}$ are linearly
disjoint over $\Q$. The Galois group of the Galois closure of their
compositum is $S_4\times\Z/2$, of order 48.
\end{theorem}

\begin{proposition}[Linear disjointness of $K_{m003}$ and $K_{m006}$]\label{prop:d6}
The cusp fields $K_{\mfld{m003}}$ and $K_{\mfld{m006}}$ are linearly
disjoint over $\Q$, with compositum degree 6. The Galois group of
the Galois closure of their compositum is $D_6$, of order 12.
\end{proposition}

Section~\ref{sec:ckm-negative} reports a negative result: no
analogous dual filling exists for the companion manifold
$M_{\mathrm{CKM}}$ within a bounded search. We regard the absence
of a positive result here as itself informative, and report it
with the same rigor as the positive results above.

\subsection*{Notation}

We adopt the following notation throughout:
\begin{itemize}
\item $M_{\mathrm{PMNS}} = \mfld{m003}(-2,3) \cong \mfld{m019}(2,1)$
      is the lepton mixing manifold.
\item $M_{\mathrm{CKM}} = \mfld{m006}(-5,2)$ is the quark mixing
      manifold.
\item $K_M$ denotes the invariant trace field of a manifold $M$,
      equivalently its cusp field for one-cusped manifolds.
\item $\mfld{m003}, \mfld{m006}, \mfld{m019}$ are one-cusped
      hyperbolic manifolds from the SnapPy census.
\end{itemize}

\section{The Two Cusped Parents}

All computations use SnapPy \cite{SnapPy} at 300-bit precision.
Polynomial identifications use \texttt{algdep} at the same
precision, with residuals below $10^{-85}$.

\begin{proposition}\label{prop:m003}
The cusp shape of \mfld{m003} is $\tau=e^{i\pi/3}$, satisfying
$x^2-x+1=0$, $\disc=-3$, $\Gal(K_{\mfld{m003}}/\Q)\cong\Z/2$.
The volume of the filled manifold $\mfld{m003}(-2,3)$ is
\[
\vol(M_{\mathrm{PMNS}}) = 0.981368828892232\ldots
\]
\end{proposition}

\begin{proof}
$e^{i\pi/3}$ is a primitive 6th root of unity, a root of the
cyclotomic polynomial $\Phi_6(x)=x^2-x+1$. The splitting field is
$\Q(\sqrt{-3})$, with Galois group $\Z/2$ over $\Q$. The volume
is computed directly by SnapPy.
\end{proof}

\begin{proposition}\label{prop:m019}
The cusp shape of \mfld{m019} satisfies the irreducible quartic
$x^4-x-1=0$, $\disc=-283$, $\Gal(K_{\mfld{m019}}/\Q)\cong S_4$.
The filling $\mfld{m019}(2,1)$ has volume
\[
\vol(\mfld{m019}(2,1)) = 0.981368828892232\ldots
\]
agreeing with $\vol(M_{\mathrm{PMNS}})$ to 15 decimal places.
\end{proposition}

\begin{proof}
$x^4-x-1$ has no rational roots and is irreducible over $\Q$
(verified by \texttt{factor} in Sage). Its discriminant is $-283$,
a negative prime. The Galois group, computed directly by Sage's
\texttt{galois\_group()}, is the transitive permutation group of
degree 4 and order 24 -- i.e.\ $S_4$ (GAP small-group id
$[24,12]$). The volume is computed directly by SnapPy.
\end{proof}

Both cusp shapes and their minimal polynomials are cross-checked
against the LMFDB number field database~\cite{LMFDB}.

\begin{remark}
The volume of the Meyerhoff manifold $M_{\mathrm{PMNS}}$ is
independently known in the literature \cite{GMM}; the value above
is consistent with the published value to the stated precision.
\end{remark}

\section{The Dual Surgery Identity}

\begin{proof}[Proof of Theorem~\ref{thm:dual}]
SnapPy computes:
\[
\vol(\mfld{m003}(-2,3)) = \vol(\mfld{m019}(2,1)) =
0.981368828892232\ldots
\]
agreeing to 15 significant figures. Both fillings have
$H_1=\Z/5$. SnapPy's \texttt{is\_isometric\_to} confirms isometry
between the two filled manifolds directly, not merely volume
coincidence. The search identifying \mfld{m019}(2,1) as a second
filling description of $M_{\mathrm{PMNS}}$ is reproduced in
\texttt{reproduce/dual\_surgery\_exploration.py}.
\end{proof}

\begin{remark}
This is a direct identity check between two specific, previously
fixed objects, not a search over candidates: \mfld{m003}(-2,3) and
\mfld{m019}(2,1) were each independently constructed before the
comparison was made.
\end{remark}

\begin{figure}[htbp]
\centering
\begin{tikzcd}
\mfld{m003}(-2,3) \arrow[dr, "\cong"] & & \mfld{m019}(2,1) \arrow[dl, "\cong"] \\
& M_{\mathrm{PMNS}} &
\end{tikzcd}
\caption{The dual surgery identity: two distinct Dehn fillings
produce the same closed manifold $M_{\mathrm{PMNS}}$.}
\label{fig:dual}
\end{figure}

\section{Linear Disjointness and Galois Closure}\label{sec:disjoint}

\begin{proof}[Proof of Theorem~\ref{thm:disjoint}]
Let $K_1=\Q(\sqrt{-3})$ (Proposition~\ref{prop:m003}, degree 2) and
$K_2=\Q[x]/(x^4-x-1)$ (Proposition~\ref{prop:m019}, degree 4). The
compositum $K_1K_2$, computed directly in Sage via
\texttt{composite\_fields}, has degree 8 over $\Q$. Since
$8 = 2\times 4 = [K_1:\Q]\cdot[K_2:\Q]$, the two fields are linearly
disjoint over $\Q$.

The compositum $K_1K_2$ is not itself Galois over $\Q$. Its Galois
closure $L$ has degree 48. Computing $\Gal(L/\Q)$ directly in Sage
gives a group of GAP small-group id $[48,48]$. An independently
constructed group, $S_4\times\Z/2$ (via \texttt{SymmetricGroup(4)}
direct product with \texttt{CyclicPermutationGroup(2)}), has the
same GAP small-group id $[48,48]$, confirming
$\Gal(L/\Q)\cong S_4\times\Z/2$.
\end{proof}

\begin{corollary}\label{cor:group}
$\Gal(L/\Q)\cong S_4\times\Z/2$, order 48.
\end{corollary}

\begin{figure}[htbp]
\centering
\begin{tikzcd}
& L & \\
& K_1K_2 \ar[dash, u, "\Gal=S_4\times\Z/2"] & \\
K_1 \ar[dash, ur, "\deg 2"] & & K_2 \ar[dash, ul, "\deg 4"'] \\
& \Q \ar[dash, ul, "\deg 2"] \ar[dash, ur, "\deg 4"'] &
\end{tikzcd}
\caption{Field tower for $K_1=\Q(\sqrt{-3})$ and $K_2=\Q[x]/(x^4-x-1)$.
The compositum $K_1K_2$ has degree 8; its Galois closure $L$ has
degree 48 with Galois group $S_4\times\Z/2$.}
\label{fig:tower}
\end{figure}

\begin{remark}[Observation only]
$S_4\times\Z/2$ is abstractly isomorphic to
$\Weyl(\SU(4))\times\Weyl(\SU(2))$, since $\Weyl(\SU(n+1))\cong
S_{n+1}$ for the $A_n$ root system and $\Weyl(\SU(2))\cong\Z/2$.
We record this as a group-theoretic observation. We do not assert,
and this paper does not attempt to establish, any physical or
field-theoretic mechanism by which the arithmetic of these cusp
fields would give rise to a gauge symmetry. That question is
explicitly open; see Section~\ref{sec:discussion}.
\end{remark}

\subsection{Linear disjointness with $K_{m006}$}

The cusp field of \mfld{m006}, the parent of $M_{\mathrm{CKM}}$,
satisfies:
\[
K_{\mfld{m006}} = \Q[x]/(x^3 + 2x + 1),\quad
\Gal(K_{\mfld{m006}}/\Q) \cong S_3,
\]
verified by the same methods as Propositions~\ref{prop:m003}
and~\ref{prop:m019}. The following result is proved analogously
to Theorem~\ref{thm:disjoint}.

\begin{proof}[Proof of Proposition~\ref{prop:d6}]
Let $K_1 = \Q(\sqrt{-3})$ (degree 2) and $K_3 = K_{\mfld{m006}}$
(degree 3). Sage computes the compositum degree as 6; since
$6 = 2\times 3$, the fields are linearly disjoint. The Galois
closure of the compositum has degree 12 and Galois group $D_6$,
the dihedral group of order 12, confirmed by Sage's
\texttt{galois\_group()} returning GAP small-group id $[12,4]$.
\end{proof}

\begin{remark}
Since $D_6$ is not a direct product of $\Z/2$ and $S_3$, this
result does \emph{not} extend the Pati--Salam correspondence of
Corollary~\ref{cor:group} to a third gauge factor. We record it
here as a proved, honestly-scoped result that does not fit the
physical narrative of this paper, rather than omit it.
\end{remark}

\section{The CKM Negative Result}\label{sec:ckm-negative}

The lepton manifold $M_{\mathrm{PMNS}}$ has a companion manifold in
the HFG program, $M_{\mathrm{CKM}}=\mfld{m006}(-5,2)$, encoding the
quark mixing matrix. Given the dual filling result of
Theorem~\ref{thm:dual}, it is natural to ask whether
$M_{\mathrm{CKM}}$ admits an analogous second filling description
from a cusped parent with $S_4$-type cusp field.

\begin{proposition}[Negative result]\label{prop:ckm-negative}
Among 9 named cusped manifolds with $S_4$ cusp field type, listed
in Table~\ref{tab:s4-manifolds}, no Dehn filling at coprime slopes
$|p|,|q|\le 15$ produces a manifold matching the volume and first
homology of $M_{\mathrm{CKM}}$.
\end{proposition}

\begin{table}[htbp]
\centering
\begin{tabular}{lccc}
\toprule
Manifold & Cusp Field Polynomial & Discriminant & Galois Group \\
\midrule
\mfld{m011} & $x^4 + 2x^3 + x^2 + x - 1$ & $-751$ & $S_4$ \\
\mfld{m019} & $x^4 - x - 1$ & $-283$ & $S_4$ \\
\mfld{m026} & $x^4 - x^3 + x^2 - x - 1$ & $-563$ & $S_4$ \\
\mfld{m033} & $x^4 - x^3 + 3x^2 - 2x + 1$ & $257$ & $S_4$ \\
\mfld{m079} & $x^4 + x^3 + 4x^2 + 3x + 2$ & $1929$ & $S_4$ \\
\mfld{m115} & $x^4 + x^3 + 2x^2 + 2x - 5$ & $-52083$ & $S_4$ \\
\mfld{m141} & $x^4 - x^3 + x - 2$ & $-1399$ & $S_4$ \\
\mfld{m155} & $x^4 - 3x^3 + 4x^2 - 6x + 5$ & $-331$ & $S_4$ \\
\mfld{m178} & $x^4 - x - 1$ & $-283$ & $S_4$ \\
\bottomrule
\end{tabular}
\caption{The nine one-cusped manifolds with $S_4$-type cusp fields
considered in the search for a dual parent of $M_{\mathrm{CKM}}$.
Polynomials and discriminants computed directly via SnapPy cusp
shapes (300-bit precision) and Sage's \texttt{algdep}, verified
against \texttt{galois\_group()} returning order 24 for all nine.
\mfld{m178} shares the same cusp field as \mfld{m019}.}
\label{tab:s4-manifolds}
\end{table}

\begin{proof}
Direct computational search, implemented in
\texttt{reproduce/dual\_surgery\_exploration.py}: for each of the 9
manifolds and each coprime slope pair $(p,q)$ with
$-15\le p,q\le 15$, $(p,q)\ne(0,0)$, the filled volume is compared
against $\vol(M_{\mathrm{CKM}})$ to a tolerance of
$10^{-6}$, restricted to fillings with $|H_1|=5$. No match is
found.
\end{proof}

\begin{remark}
This is a bounded search, not an exhaustive census; it does not
rule out a dual parent at larger slopes, among unlisted manifolds,
or of a different cusp-field type. We report it as a negative
result within the stated bounds, not as a proof of non-existence.
Two smaller bounds ($\le 12$, $\le 20$) have circulated elsewhere
in the corpus; neither traces to this script, and the figure
$|p|,|q|\le 15$ stated here is the one actually verified against
the code.
\end{remark}

\section{Discussion}\label{sec:discussion}

The results above are, we believe, a complete and self-contained
mathematical statement: a proved identity (Theorem~\ref{thm:dual}),
a proved arithmetic independence result with an explicitly computed
Galois group (Theorem~\ref{thm:disjoint}), a proved linear
disjointness result for a second field pair (Proposition~\ref{prop:d6}),
and a clearly bounded negative result (Proposition~\ref{prop:ckm-negative}).
Three questions remain explicitly open and are not addressed by this
paper:

\begin{itemize}
\item \textbf{The $\SU(2)_R$ factor.} The group
$S_4\times\Z/2\cong\Weyl(\SU(4))\times\Weyl(\SU(2))$ matches only
the left-handed factor of the Pati--Salam gauge group
$\SU(4)\times\SU(2)_L\times\SU(2)_R$. No cusped parent or geometric
construction addressing an $\SU(2)_R$ factor is known.

\item \textbf{The physical mechanism.} No field-theoretic
derivation connects the abstract group isomorphism of
Corollary~\ref{cor:group} to an actual gauge theory. Whether a
construction along the lines of the Dimofte--Gaiotto--Gukov 3d-3d
correspondence \cite{DGG} could supply such a mechanism is an
open research question, not addressed here.

\item \textbf{Linear disjointness with a third field.}
Proposition~\ref{prop:d6} shows that $K_{\mfld{m003}}$ and
$K_{\mfld{m006}}$ are also linearly disjoint over $\Q$, with
Galois closure $D_6$ (dihedral, order 12) -- not $S_3\times\Z/2$.
Because $D_6$ is not a direct product of the individual Weyl
groups $\Z/2$ and $S_3$, this result does \emph{not} extend the
Pati--Salam correspondence of Corollary~\ref{cor:group} to a third
gauge factor. We record it as a proved, honestly-scoped result
that does not fit the physical narrative of this paper, rather
than omit it.
\end{itemize}

No claim in this paper extends beyond what is directly verified by
SnapPy or Sage computation. The group-theoretic observation of
Section~\ref{sec:disjoint} is explicitly not a physical derivation.

\appendix
\section{Computational Reproducibility and Certification}\label{app:repro}

This appendix records the computational provenance of the results above.
The purpose is to distinguish high-precision numerical evidence from
rigorously certified algebraic or topological statements.

\begin{enumerate}
\item \textbf{Census objects and fillings.}
The manifolds are identified by their SnapPy census names and the Dehn
filling slopes are given in SnapPy's peripheral-curve convention.
The exact software version and census version used for the archived
computation should be recorded with the accompanying source release.

\item \textbf{Hyperbolic computations.}
Volumes and cusp shapes are numerical outputs. The working precision
used in the computations reported here is 300 bits. Decimal agreement
is used only as a consistency check and is not, by itself, taken as a
proof of isometry.

\item \textbf{Dual-filling identification.}
The decisive computational step is the direct isometry test between
the two filled manifolds. The archived reproduction script records
the precise SnapPy routine used. For archival publication, the
authors should additionally retain any available certified
triangulation or combinatorial isomorphism certificate produced by
the software, so that the topological identification does not rest
on equality of numerical invariants alone.

\item \textbf{Number fields.}
Minimal polynomials inferred from high-precision cusp data are recorded
together with the reconstruction procedure and residual. Irreducibility,
discriminants, composita, normal closures, and finite Galois groups are
then computed exactly in SageMath once the defining polynomial has been
fixed. A numerical algebraic-dependence calculation is therefore
distinguished from the subsequent exact algebraic computations.

\item \textbf{Archival record.}
The reproduction repository should be cited by immutable commit hash
and, for submission, preferably archived with a DOI-bearing release.
The archive should include software versions, commands, raw output, and
the scripts used to generate every computational statement in the paper.
\end{enumerate}

No physical interpretation is required for any theorem in this paper.
The notation $M_{\mathrm{PMNS}}$, where mentioned, records the terminology
of the broader HFG programme and is not part of the mathematical content
of the dual-filling or Galois-group results.

\section*{Data availability}
Scripts at
\url{https://github.com/drmlgentry/hyperbolic-flavor-geometry}
(\url{reproduce/dual_surgery_exploration.py}).

\section*{Acknowledgements}
During preparation of this manuscript the author used Claude as a
research assistant. All mathematical results were independently
verified using SnapPy and SageMath. The author takes full
responsibility for all content.

\begin{thebibliography}{99}

\bibitem{GMM}
D.~Gabai, R.~Meyerhoff, P.~Milley,
J.\ Amer.\ Math.\ Soc.\ \textbf{22}, 1157 (2009).

\bibitem{PS}
J.~C.~Pati and A.~Salam,
Phys.\ Rev.\ D \textbf{10}, 275 (1974).

\bibitem{SnapPy}
M.~Culler, N.~M.~Dunfield, M.~Goerner, J.~R.~Weeks,
SnapPy, \url{http://snappy.computop.org}.

\bibitem{LMFDB}
The LMFDB Collaboration,
\emph{The L-functions and Modular Forms Database},
\url{http://www.lmfdb.org}.

\bibitem{DGG}
T.~Dimofte, D.~Gaiotto, S.~Gukov,
``Gauge Theories Labelled by Three-Manifolds,''
Commun.\ Math.\ Phys.\ \textbf{325}, 367 (2014),
arXiv:1108.4389.

\end{thebibliography}

\end{document}
